% F081: precise mathematical teaching diagram.
\begin{tikzpicture}[x=1cm,y=1cm,>=stealth,font=\small]
\draw[rounded corners=3pt,gray!40] (0,0) rectangle (14.2,9.0);
\begin{scope}[yshift=1cm]
\node at (7.1,7.4) {内部最高点：两侧割线逼向同一条水平切线};
\begin{scope}[shift={(4.2,2.5)},x=2.3cm,y=2cm]
\draw[->,gray!65] (-1.35,0)--(1.35,0) node[right] {$x$};
\draw[->,gray!65] (0,-0.3)--(0,1.45) node[above] {$y$};
\draw[blue!75!black,line width=1.15pt,,domain=-1:1,samples=160] plot (\x,{1-(\x)*(\x)});
\fill[blue!75!black] (-1,0) circle (1.8pt);\node[below left] at (-1,0) {$A$};
\fill[blue!75!black] (1,0) circle (1.8pt);\node[below right] at (1,0) {$B$};
\fill[blue!75!black] (0,1) circle (1.8pt);\node[above left] at (0,1) {$C$};
\draw[orange!85!black,line width=1.2pt] (-0.35,1)--(0.35,1);
\draw[green!50!black,line width=1pt] (-.7,.51)--(0,1);\draw[purple,line width=1pt] (0,1)--(.7,.51);
\fill[green!50!black] (-0.7,0.51) circle (1.8pt);\node[above left] at (-0.7,0.51) {$h<0$};
\fill[purple] (0.7,0.51) circle (1.8pt);\node[above right] at (0.7,0.51) {$h>0$};
\draw[dashed,gray!60] (0,0)--(0,1);\node[below right] at (0,0) {$c$};
\end{scope}
\node[green!50!black] at (10.4,5.3) {左侧差商 $\ge0$};
\node[purple] at (10.4,4.3) {右侧差商 $\le0$};
\node[] at (10.4,3.2) {可导：两侧极限相同};
\node[] at (10.4,2.1) {所以 $f^{\prime}(c)=0$};
\end{scope}
\node[text width=13cm,align=center] at (7.1,.9) {闭区间连续保证能取到最高值；两端等高且有更高点，最高点就在内部。};
\end{tikzpicture}
